\section{Realizacija}
\label{sec:realizacija}

Planas realizuotas TypeScript kalba, naudojant \ref{subsec:biblioteka} poskyryje išvardytas bibliotekas. Programa padalyta į modulius pagal atsakomybes (\ref{tab:moduliai} lentelė): pagrindinis failas \texttt{main.ts}, o kiti moduliai -- kataloge \texttt{code}. Programa paleidžiama komanda \texttt{npx tsx 5lab/main.ts}.

\begin{table}[H]
    \centering
    \caption{Programos moduliai.}
    \label{tab:moduliai}
    \begin{tabularx}{\textwidth}{l L}
        \toprule
        \textbf{Failas} & \textbf{Paskirtis} \\
        \midrule
        \texttt{code/interfaces.ts} & Vartotojo įrašo tipas \texttt{User} (\ref{tab:db} lentelės laukai ir \texttt{passwordHash}) \\
        \texttt{code/users-db.ts} & Vartotojų „duomenų bazė“ ir registracija (\texttt{register}) \\
        \texttt{code/secret-encryption.ts} & Paslapties šifravimas ir iššifravimas (\texttt{encrypt}, \texttt{decrypt}) \\
        \texttt{code/totp-verification.ts} & TOTP kodo tikrinimas (\texttt{acceptedStep}) \\
        \texttt{code/mfa-activation.ts} & MFA aktyvavimas (\texttt{mfaSetup}, \texttt{mfaConfirm}) \\
        \texttt{code/login.ts} & Prisijungimas ir blokavimas (\texttt{login}, \texttt{recordFailure}) \\
        \texttt{main.ts} & Demonstracinis scenarijus, kurio išvestis pateikta \ref{subsec:rezultatai} poskyryje \\
        \bottomrule
    \end{tabularx}
\end{table}

Palyginti su planu, realizacija supaprastinta:
\begin{itemize}
    \item vietoj duomenų bazės naudojama \texttt{Map} struktūra, kurios raktas -- vartotojo vardas, o reikšmė -- \texttt{passwordHash} ir \ref{tab:db} lentelės laukai;
    \item nėra HTTP sąsajos ir sesijų, todėl abu prisijungimo žingsniai sujungti į vieną funkciją \texttt{login};
    \item šifravimo raktas sugeneruojamas paleidus programą, o el. laiškai nesiunčiami.
\end{itemize}

Paslaptis šifruojama ir iššifruojama pagalbinėmis funkcijomis \texttt{encrypt} ir \texttt{decrypt}: naudojamas AES-256-GCM su atsitiktiniu 96 bitų IV, o duomenų bazėje saugomi IV, autentifikavimo žyma ir šifrograma, atskirti dvitaškiais.

\subsection{MFA aktyvavimas}
\label{subsec:realizacija_aktyvavimas}

Funkcija \texttt{mfaSetup} sugeneruoja paslaptį, išsaugo ją užšifruotą ir sukuria QR kodą. Funkcija \texttt{mfaConfirm} įjungia MFA tik tada, kai vartotojas įveda teisingą pirmą kodą.

\begin{lstlisting}
export async function mfaSetup(username: string): Promise<string> {
  const user = users.get(username)!;
  if (user.mfaEnabled) {
    throw new Error('MFA is already enabled');
  }

  const totp = new TOTP({
    issuer: ISSUER,
    label: username,
    secret: new Secret({ size: 20 }),
  });
  // Saved now, but MFA stays off until the user confirms the first code.
  user.totpSecret = encrypt(totp.secret.base32);

  const uri = totp.toString();
  await QRCode.toFile(QR_CODE_FILE, uri, { scale: 8 });
  return uri;
}

export function mfaConfirm(username: string, code: string): boolean {
  const user = users.get(username)!;
  if (user.mfaEnabled) {
    return false;
  }

  const step = acceptedStep(user, code);
  if (step === null) {
    return false;
  }
  user.mfaEnabled = true;
  user.lastUsedStep = step;
  return true;
}
\end{lstlisting}

\subsection{Autentifikavimas}
\label{subsec:realizacija_autentifikavimas}

Kodą tikrina funkcija \texttt{acceptedStep}, kurią naudoja ir aktyvavimas, ir prisijungimas. Ji grąžina priimto kodo laiko intervalą arba \texttt{null}, jei kodas neteisingas arba jau panaudotas.

\begin{lstlisting}
export function acceptedStep(user: User, code: string): number | null {
  if (user.totpSecret === null) {
    return null;
  }

  const totp = new TOTP({ secret: decrypt(user.totpSecret) });
  const now = Date.now();
  const delta = totp.validate({ token: code, timestamp: now, window: 1 });
  if (delta === null) {
    return null;
  }

  const step = totp.counter({ timestamp: now }) + delta;
  return step > user.lastUsedStep ? step : null;
}
\end{lstlisting}

Funkcija \texttt{login} pirmiausia patikrina, ar paskyra neužblokuota, tada slaptažodį ir tik po to TOTP kodą. Kiekvieną nesėkmę užfiksuoja funkcija \texttt{recordFailure}.

\begin{lstlisting}
export async function login(
  username: string,
  password: string,
  code: string
): Promise<boolean> {
  const user = users.get(username);
  if (!user || Date.now() < user.lockedUntil) {
    return false;
  }

  // Factor 1: the password, checked against its Argon2id hash.
  if (!(await argon2.verify(user.passwordHash, password))) {
    recordFailure(user);
    return false;
  }

  // Factor 2: the server computes the expected codes itself and compares.
  if (user.mfaEnabled) {
    const step = acceptedStep(user, code);
    if (step === null) {
      recordFailure(user);
      return false;
    }
    user.lastUsedStep = step;
  }

  user.failedAttempts = 0;
  return true;
}

function recordFailure(user: User): void {
  user.failedAttempts++;
  if (user.failedAttempts >= MAX_FAILED_ATTEMPTS) {
    user.failedAttempts = 0;
    user.lockedUntil = Date.now() + LOCK_DURATION_MS;
  }
}
\end{lstlisting}

\subsection{Rezultatai}
\label{subsec:rezultatai}

Programa užregistruoja vartotoją, aktyvuoja jam MFA ir išbando kelis prisijungimo atvejus. Autentifikatoriaus programėlę imituoja \texttt{TOTP} objektas, sukurtas vien iš QR kode esančio \texttt{otpauth://} adreso. Programos išvestis:

\begin{Verbatim}[frame=single, fontsize=\footnotesize, formatcom=\setstretch{1}]
Secret for manual entry: 2QQO6FQ7ZNQHMVVV5XBAAFLH5ME23YJD
QR code saved: totp-qr.png
DB field totpSecret: 3chGKc/fG9GG1sAS:cm++vQcJdcTftoHnS4w1uA=...
Activation with wrong code 123456 -> rejected
Activation with code 810340 -> accepted
Login with wrong password and code 025236 -> rejected
Login with already used code 810340 -> rejected
Login with next-step code 025236 -> accepted
5 wrong codes in a row -> account locked until 21:21:40
\end{Verbatim}

\ref{fig:qr} pav. pateiktame QR kode užkoduotas \texttt{otpauth://} adresas su ta pačia paslaptimi, kuri išvestyje parodyta rankiniam įvedimui.

\begin{figure}[H]
    \centering
    \includegraphics[width=4cm]{images/totp-qr.png}
    \caption{Programos sugeneruotas aktyvavimo QR kodas.}
    \label{fig:qr}
\end{figure}

Rezultatai atitinka planą:
\begin{itemize}
    \item duomenų bazėje paslaptis saugoma užšifruota, o ne Base32 pavidalu;
    \item neteisingas pirmasis kodas MFA neįjungia;
    \item su neteisingu slaptažodžiu prisijungti nepavyksta net su tinkamu kodu -- tas pats kodas po to priimamas su teisingu slaptažodžiu;
    \item aktyvavimo kodas antrą kartą nepriimamas, nors pagal laiką jis dar tinkamas;
    \item kodas, kurį programėlė rodytų tik po 30 sekundžių, priimamas dėl leidžiamo vieno intervalo nuokrypio;
    \item po 5 neteisingų kodų iš eilės paskyra užblokuojama 15 minučių.
\end{itemize}
